% sections/03_regime_map.tex -- Map of regimes (Q5, hypothesis H4). SKELETON.
% Source: PLAN.md Q5; brief notes/briefs/regime-map.md; results in experiments/regime_map/ and
% notes/regime_map.md; theory in notes/theory_regimes.md. Labels follow tab:regimes exactly.
\section{Map of regimes}\label{sec:regime-map}

\placeholder{Skeleton. Phase diagrams of the training regimes over initialization scale, bond
dimension, number of cores and core imbalance, with predicted boundaries and threshold
sensitivity. Waits for the regime-map results.}
% TODO(status): PLAN.md tags Q5 as "known in DLNs/diagonal nets; open for TTNs". Mark the TTN
% phase diagrams \newresult, and the reproduction of Chizat/Woodworth/Kunin effects \knownresult.

\subsection{Protocol}\label{sec:regime-protocol}
\placeholder{Centred model as the one-knob family; labels from \cref{tab:regimes}; seeds, sweeps,
time budgets.}
% TODO: centred model (target W* + W_0, Sec. 2.4 variant (b)) as the main family; two-copy model
% with mask as cross-check; large sigma needs time rescaled by the kernel scale sigma^{2(K-1)}.
% Per grid point >= 5 seeds: delta_Theta (P = 128 probes), delta_theta, separation ratios, label,
% final loss. Report label fractions per point, not only the majority label.

\subsection{Initialization scale and bond dimension}\label{sec:regime-sigma-chi}
\placeholder{Phase diagram $\isc\times\chi$ for binary trees with $\Ns=4,8$.}
% TODO: sigma over >= 3 decades (deep S2S to lazy), chi from the teacher rank to >= 8x it; odeco
% teacher with 3 components, weights [3, 2, 1]. Figures: categorical phase diagram with predicted
% boundaries overlaid; continuous heatmaps of delta_Theta(T) and min_i rho_i; representative runs of
% each regime. Theory: lazy boundary as a threshold on the relative kernel change (Geiger et al.;
% Chizat et al.'s relative scale kappa, Sec. 1.2, compares the relative change of Dh with the relative
% loss decrease, not weight vs function change); S2S boundary from the escape-time formulas with
% separation ratio rho = 3.

\subsection{Number of cores and function-space scale}\label{sec:regime-K}
\placeholder{Phase diagram $\isc\times\Kc$; boundaries in both $\isc$ and $\fsc=\isc^{\Kc}$.}
% TODO: binary N = 4, 8, 16 and MPS N = 4..12 (K from 3 to ~15). Test the prediction that the S2S
% boundary sigma_c(K) moves toward 1 as K grows (H4): plot vs sigma and vs alpha = sigma^K.

\subsection{Imbalance between cores}\label{sec:regime-imbalance}
\placeholder{Does the conserved imbalance $\Q{e}$ move the TTN between regimes, as in deep
linear networks?}
% TODO: per-core scales (root vs rest) at fixed function-space scale prod_v sigma_v; compare with
% Kunin et al. 2024 (kunin2024get) and Domine et al. (lazy-to-rich exact solutions).

\subsection{Large bond dimension}\label{sec:regime-large-chi}
\placeholder{NTK-type versus $\mu$P-type parameterizations as $\chi$ grows.}
% TODO: chi = 4..64; per-core init scales and lr multipliers from notes/theory_regimes.md (spectral
% condition of Yang, Simon & Bernstein 2023 applied per core); delta_Theta and delta_theta vs chi;
% does feature learning survive? Status: \derived for the scalings, \empirical for the curves.

\subsection{Threshold sensitivity and cross-checks}\label{sec:regime-sensitivity}
\placeholder{How the boundaries move with the kernel threshold $\{0.05,0.1,0.2\}$ and the
separation threshold $\{2,3,5\}$; two-copy cross-check.}
% TODO: one panel per threshold pair; state which boundaries are robust. Two-copy model at a few
% grid points must give the same labels as the centred model (or explain why not).

\subsection{Summary: which parameter moves which boundary}\label{sec:regime-summary}
\placeholder{Table: parameter, boundary moved, direction, fitted versus predicted.}
